\documentclass[11pt]{article}
\usepackage[a4paper,margin=19mm]{geometry}
\usepackage{fontspec}
\setmainfont{lmroman10-regular.otf}
\usepackage{amsmath,amssymb,amsthm,mathtools,xcolor,graphicx,hyperref}
\hypersetup{hidelinks}
\numberwithin{equation}{section}
\newtheorem{theorem}{Theorem}[section]
\newtheorem{lemma}[theorem]{Lemma}
\newtheorem{proposition}[theorem]{Proposition}
\newtheorem{corollary}[theorem]{Corollary}
\newtheorem{claim}[theorem]{Claim}
\newtheorem{fact}[theorem]{Fact}
\newtheorem{conjecture}[theorem]{Conjecture}
\theoremstyle{definition}
\newtheorem{definition}[theorem]{Definition}
\newtheorem{example}[theorem]{Example}
\newtheorem{remark}[theorem]{Remark}
\newtheorem{question}[theorem]{Question}
\theoremstyle{plain}
\newcommand{\C}{\mathbb C}
\newcommand{\R}{\mathbb R}
\newcommand{\Q}{\mathbb Q}
\newcommand{\Z}{\mathbb Z}
\newcommand{\N}{\mathbb N}
\newcommand{\eee}{\mathrm e}
\newcommand{\cclass}[1]{\textsf{#1}}
\newcommand{\cproblem}[1]{\textsc{#1}}
\newcommand{\ie}{i.\,e.}
\newcommand{\Ie}{I.\,e.}
\newcommand{\eg}{e.\,g.}
\newcommand{\Eg}{E.\,g.}
\newcommand{\phd}{Ph.\,D.}
\newcommand{\msc}{M.\,S.}
\newcommand{\bsc}{B.\,S.}
\newcommand{\expref}[2]{\hyperref[#2]{#1~\ref{#2}}}
\newcommand{\expeqref}[2]{\hyperref[#2]{#1~\eqref{#2}}}
\newcommand{\exprefsilent}[2]{\hyperref[#2]{#1}}
\usepackage{eprint}
\DeclareMathOperator{\Supp}{supp}
\DeclareMathOperator{\spa}{sparsity}
\DeclareMathOperator{\dimn}{dim}
\DeclareMathOperator{\vspan}{span}
\DeclareMathOperator{\Add}{Add}   %

\newcommand{\etal}{{\it et~al.}}
\newcommand{\lref}[2][]{\hyperref[#2]{#1~\ref*{#2}}}
\newcommand{\F}{{\mathbb F}}

\newtheorem{clm}[theorem]{Claim}
\newtheorem{observation}[theorem]{Observation}
\usepackage{comment}

\title{Near-square-root Fourier dimension bound from Boolean Fourier sparsity}
\author{Swagato Sanyal}
\date{doi:10.4086/toc.2019.v015a011; CC BY 3.0}
\begin{document}
\maketitle
\noindent\textit{Editorial note.} Full final author text and proof retained. Earlier weaker arXiv bound not counted separately. Discussion conjectures remain conjectures.\par\medskip
\noindent\textit{Allowed standard prerequisites (editorial).}\begin{itemize}
\item Boolean Fourier expansion, uncertainty principle and affine monochromatic restriction theorem. Original standard prerequisite statements and citations retained; local non-adaptive algorithm and support argument retained in full.. No recursive reproof of cited standard results.
\end{itemize}\medskip
\section{Introduction}
\label{intro}
The study of Boolean functions involves studying various complexity measures and their inter-relation\-ships. Two such measures,
which we investigate in this article, are the Fourier dimension and
the Fourier sparsity. Let
$f:\mathbb{F}_2^n \to \{1,-1\}$ be a Boolean function with Fourier expansion
\[f(x)=\sum_{\gamma \in \mathbb{F}_2^n}\widehat{f}(\gamma) \chi_\gamma(x).\]
where $\chi_\gamma(x):=(-1)^{\sum_{i=1}^n \gamma_ix_i}$ are the characters and the $\widehat{f}(\gamma)$'s are real numbers, called the Fourier coefficients of $f$.
The Fourier dimension and Fourier sparsity are defined as follows.
\begin{definition}[Fourier sparsity and dimension]
 For a Boolean function $f:\mathbb{F}_2^n \to \{1,-1\}$ with Fourier expansion
\[f(x)=\sum_{\gamma \in \mathbb{F}_2^n}\widehat{f}(\gamma)
\chi_\gamma(x),\] the \emph{Fourier support} of $f$, denoted by $\Supp(\widehat{f})$,
is defined as 
\[\Supp(\widehat{f}) := \{\gamma \in \mathbb{F}_2^n: \widehat{f}(\gamma) \neq 0\}.\]
The \emph{Fourier sparsity} of $f$, denoted by $\spa(\widehat{f})$, is defined as the size of the Fourier support of $f$,
\ie, 
\[\spa(\widehat{f}) :=|\Supp(\widehat{f})|,\]
while the \emph{Fourier dimension} $\dimn(\widehat{f})$ of $f$ is defined as the
dimension of the span of $\Supp(\widehat{f})$.
\end{definition}
Note that the Fourier expansion of $f$ is a
%
multilinear   %
polynomial in the variables $y_i:=(-1)^{x_i}$. With respect to this view, Fourier sparsity of a Boolean function $f$ is the number of monomials that appear in the Fourier expansion of $f$. Fourier dimension, on the other hand, is the smallest number of monomials, or equivalently parity functions, whose values always determine $f$. It is natural to investigate the power and limitation of polynomials with low Fourier sparsity or Fourier dimension. Fourier sparsity and Fourier dimension were studied by Gopalan \etal~\cite{DBLP:journals/siamcomp/GopalanOSSW11} in the context of property testing. More recently these quantities have been studied in the context of learning~\cite{DBLP:journals/toct/HavivR16, DBLP:journals/corr/abs-1810-00481}. An approximate analog of Fourier dimension has been shown to characterize the randomized one-way communication complexity of an important and well-studied subclass of functions called the XOR functions, over uniformly distributed inputs~\cite{DBLP:conf/coco/KannanMSY18}. A Boolean function $F(x,y)$ on two $n$-bit inputs is an XOR function if there exists a Boolean function $f$ on $n$ bits such that
$F(x,y)=f(x \oplus y)$. %

Besides, the study on Fourier sparsity has attracted attention of complexity theorists due to its intimate connection to the \textit{log-rank conjecture} of communication complexity. Fourier dimension is related to a simultaneous communication game. This connection is discussed in more detail later in this section.

The following inequalities easily follow from the definition of Fourier sparsity and dimension%
.
\begin{equation}
\label{eq:order}
 \log_2 \spa(\widehat{f}) \leq \dimn(\widehat{f}) \leq \spa(\widehat{f}).
\end{equation}
There are functions (e.g., indicator functions of subspaces) for which
the first inequality is tight (\ie, %
holds with equality). In this
%
note   %
we examine the tightness of the second inequality. Note that the Fourier
transform of a Boolean function is a polynomial with a very special
%
property;     %
it evaluates to $1$ or $-1$ on all inputs from $\{1,-1\}^n$. It is thus natural to expect that there is always a good amount of dependency amongst its monomials. The gap between Fourier sparsity and Fourier dimension is one way of quantifying this dependency.

For the second inequality, the \textit{address function} is one function
having asymptotically the closest gap between Fourier dimension and sparsity.\footnote{Recently Khalyavin, Lobanov and Tarannikov have constructed a function for which the gap between Fourier dimension and Fourier sparsity is closer by a constant factor than in the address function~\cite{khalyavin}.} Let $s$ be a power of $2$. The address function
 $\Add_s:\{0,1\}^{(1/2) \log s +
  \sqrt{s}}\to \{1,-1\}$ is defined
as $$\Add_s(x,y_1,y_2,\dots,y_{\sqrt{s}}) := (-1)^{y_x}, \quad x \in \{0,1\}^{\frac12\log s}, y_i \in \{0,1\},$$ where $x$ is interpreted as an address in $\{1, \ldots, \sqrt s\}$. In other words, on every input
$(x,y)$, $\Add_s(x,y)$ is the value of the
%
addressed    %
input bit $y_x$
indexed by the addressing variables $x$. The address function has
Fourier sparsity $s$ and Fourier dimension at least $\sqrt{s}$. To see this, note that
\[
  \Add_s(x,y)=\sum_{\tilde{x} \in \{0,1\}^{({1}/{2}) \log s}} \mathbf{1}_{x=\tilde{x}}\cdot (-1)^{y_{\tilde{x}}}.
\]
Each indicator function $\mathbf{1}_{x=\tilde{x}}$ has Fourier
%
%
sparsity\footnote{Note    %
  that the indicator functions evaluate to $\{0,1\}$ instead of $\{1,-1\}$. However, the theory of Fourier analysis extends to such functions as well.}
$\sqrt{s}$.   %
Since the summation is over $\sqrt{s}$ terms, and since there is no
%
%
cancellation due to the presence of a character corresponding to a fresh variable $y_{\tilde{x}}$ in each term, the Fourier sparsity of the address function is equal to $\sqrt s\cdot \sqrt s=s$. To see that the dimension is at least $\sqrt{s}$, note that for each $\tilde{x} \in \{0,1\}^{({1}/{2})\log s}$, the character $(-1)^{y_{\tilde{x}}}$ appears in the Fourier transform.
%
%
We prove that for any Boolean function, this is the asymptotically highest possible value of $\dimn(\widehat{f})$ in terms of
$\spa(\widehat{f})$, up to a factor of $O(\log s)$. This was presented as an open problem at the Simons workshop on Real Analysis in Testing, Learning and Inapproximability, 2013 by John Wright.

Our main result is the following. %

\begin{theorem}
 \label{thm:th1}
Let $f$ be a Boolean function with $\spa(\widehat{f})=s$. Then, $\dimn(\widehat{f})=O\left(\sqrt{s} \log s \right)$.
\end{theorem}

%
In a  %
preliminary version of this
%
article, %
%
%
posted on \href{http://arxiv.org/abs/1407.3500}{arXiv} and   %
\href{https://eccc.weizmann.ac.il/report/2014/088/}{ECCC},   %
we proved an upper bound of $O(s^{2/3})$ on $\dimn(\widehat{f})$.
Avishay Tal observed that the analysis can be tightened to
%
obtain  %
the
near-optimal upper bound of $O\left(\sqrt{s} \log s \right)$.

Prior to this work, Gavinsky \etal~\cite{alex} had proved that the dimension of any Boolean function with Fourier sparsity $s$ is $O(s/\log s)$.

\lref[Theorem]{thm:th1} is proved using a lemma of
Tsang~\etal~\cite{tsang} bounding the codimension of an affine subspace
restricted to which the function is constant, in terms of the
Fourier sparsity of the function.
%
The following result is a corollary to~\cite[Lemma~28]{tsang}.
%
\begin{lemma}[Tsang, Wong, Xie, and Zhang]
\label{lem:set}
Let $f:\mathbb{F}_2^n \rightarrow \{1,-1\}$ be a Boolean function
with Fourier sparsity $s$. Then there is an affine subspace $V$ of $\mathbb{F}_2^n$ of codimension $O(\sqrt{s})$ such
that $f$ is constant on $V$.
\end{lemma}
%
\paragraph{Proof idea of \lref[Theorem]{thm:th1}.}
%
We begin
%
with  %
a simple but crucial
observation made by Gopalan \etal~\cite{DBLP:journals/siamcomp/GopalanOSSW11}, that the Fourier dimension of a Boolean function is
equivalent to its non-adaptive parity decision tree complexity (see
\lref[Proposition]{prop:basic}). This offers us a potential approach
towards
%
an upper bound on  %
the Fourier dimension of a Boolean function:
exhibiting a shallow non-adaptive parity decision tree of the
function. We recall that the character functions essentially compute the parities of various subsets of the input bits. Thus a parity decision tree can be thought of as querying various character functions. Parity functions, in turn, are linear functions from $\F_2^n$ to $\F_2$; thus affine subspaces can be described by specifying a set of parities that are set to various values.

Towards this end, we first recall the construction of an
(adaptive) parity decision tree for a Boolean function $f$ of Fourier sparsity $s$ by Tsang~\etal~\cite{tsang}, which in
turn improves on an earlier construction due to 
Shpilka~\etal~\cite[Theorem~1.1]{DBLP:conf/innovations/ShpilkaTV14}.
The broad idea of their
construction is as follows. %
At any point in time, a partial
tree is maintained whose leaves are functions which are restrictions of $f$ on different affine subspaces. Then a non-constant
leaf $\tau$ is picked arbitrarily, and a small set
of linear restrictions is obtained by invoking \lref[Lemma]{lem:set}, such that the restricted function $f\mid_\tau$ at that leaf becomes
constant. The next step is observing that if $f\mid_\tau$ is further restricted to
all the affine subspaces obtained by setting the same set of 
parities in all possible ways, then the Fourier sparsity of each of the corresponding restricted functions
is at most half of that of $f\mid_\tau$. This is because, in
the former restriction, since the function becomes constant, the
Fourier coefficients corresponding to non-constant characters must
disappear in the restricted space. This can only happen if every
non-constant character gets identified with at least one
other character. This identification leads to halving of the Fourier
support. Note that by \lref[Lemma]{lem:set} the number of queries we have spent to achieve this reduction in Fourier sparsity is
\[
  O\left(\sqrt{\spa(\widehat{f\mid_\tau})}\right).
\]
A calculation gives us that proceeding in this way a parity decision tree
of depth $O(\sqrt{s}$) is obtained. 


Note that the choice of parities restricted at various steps depends on the leaf (function)
chosen, and hence on the outcomes of the preceding queries. Thus the
constructed tree is an adaptive one. In this article, we make this
tree non-adaptive, at the cost of a logarithmic increase in depth. At each
step, we choose an appropriate function (leaf) of the partial tree constructed thus far, invoke
\lref[Lemma]{lem:set}, and obtain restrictions which make it
constant. Then we query the same set of parities at every
leaf. Note that doing that in each step results in a non-adaptive tree, as the set of parities
queried at each step is the same for all the leaves of the current partial tree, and hence independent of the responses to the previous queries. Then we argue that this leads to a significant
reduction of Fourier sparsity. 
Let $s^{(i)}$ be the Fourier sparsity of the function (leaf)
chosen at the $i$-th step. It can be shown that, in the next step, 
the size $\ell^{(i)}$ of the union of the Fourier supports of all the leaves falls at least
by $s^{(i)}/2$. {}From \lref[Lemma]{lem:set}, the number of queries
spent in the $i$-th step is $O(\sqrt{s^{(i)}})$. Using the
\emph{Uncertainty Principle} (\lref[Theorem]{thm:up}) we show that 
$s^{(i)} \geq \left(\ell^{(i)}\right)^2/s$ (see \expref{Lemma}{lem:support}). We combine all these facts to show that continuing in
this fashion, in a small number of steps and making at most $O(\sqrt{s} \log s)$
queries, all the leaves become constant functions. The details of the construction of the
non-adaptive parity decision tree, and its analysis, are given in \lref[Section]{themain}.
\paragraph{Connections to communication complexity and the log-rank
conjecture.}  %
The log-rank conjecture is a long-standing and
important conjecture in communication complexity. The statement of the
conjecture is that the deterministic communication complexity of a
Boolean function is asymptotically bounded above by some fixed
polylogarithm of the rank (over the real numbers) of its communication
matrix. The best known upper bound of deterministic communication
complexity of a function in terms of the rank is $O(\sqrt{\mbox{rank}} \log \mbox{rank})$ due to Lovett~\cite{lovett}.\footnote{We note here that the related \emph{log approximate-rank conjecture} has
%
recently been refuted by
Chattopadhyay \etal~\cite{DBLP:journals/eccc/ChattopadhyayMS18}, and
the function
%
they use   %
is an XOR function.}
The rank of the communication matrix of an XOR 
function $F(x,y):=f(x \oplus y)$ is known to be equal to the Fourier sparsity $s$ of $f$. For the special case of XOR functions the result by Lovett also follows from
%
a result of  %
of Tsang~\etal~\cite{tsang},\footnote{For
XOR functions, the communication complexity upper bound is in fact $O(\sqrt{\mbox{rank}})$ which is better than $O(\sqrt{\mbox{rank}} \log \mbox{rank})$ by a logarithmic factor. It follows from the parity decision tree complexity upper bound proved by Tsang \etal~\cite{tsang}.}
which improves on
%
a result by   %
Shpilka~\etal~\cite{DBLP:conf/innovations/ShpilkaTV14}. A consequence of \lref[Theorem]{thm:th1} is that XOR functions admit a protocol of complexity $O(\sqrt{\mbox{rank}} \log \mbox{rank})=O(\sqrt{s} \log s)$ in the simultaneous communication model.\footnote{We thank the anonymous referee for pointing this out.} We note that both the earlier protocols~\cite{lovett, tsang} are two-way. The simultaneous protocol is as follows. Alice and Bob a priori agree on a set $\cal S$
of $O(\sqrt{s} \log s)$ parities that span the Fourier support
of $F(x,y)=f(x \oplus y)$. The existence of $\cal S$ is guaranteed by \lref[Theorem]{thm:th1}. The character corresponding to each
parity $P$ in $\cal S$ is a product of a character $M_x$ in variables in $x$ and a character $M_y$ in variables in $y$.
Upon receiving inputs $x$ and $y$, Alice and Bob compute the values of $M_x$ and $M_y$, respectively, %
for each $P \in \cal S$, and send the evaluations to the referee.
The referee then can evaluate the characters corresponding to each parity $P \in \cal S$ on the input $x \oplus y$. Since every other character
in the Fourier expansion of $F$ is determined by the characters in $\cal S$, the referee can compute the value of all those characters. Finally, the referee evaluates $F(x,y)$ from its Fourier expansion and outputs it.

\paragraph{Some remarks about \lref[Lemma]{lem:set}.}
%
\lref[Lemma]{lem:set} is not believed to be tight. 
Tsang \etal~\cite{tsang} investigated this question while studying
the log-rank conjecture for XOR functions. They suggested a direction towards proving
the log-rank conjecture for XOR
functions. In particular,
%
they  %
proposed a protocol for
an XOR function $F(x,y)=f(x \oplus y)$ based on a parity decision tree of $f$ and showed that the
communication complexity of the proposed protocol is polylogarithmic
in the rank of the communication matrix if the following related
%
%
conjecture, stated as~\cite[Conjecture~27]{tsang}, 
is true.
\begin{conjecture}[Tsang \emph{et al.}]    %
 \label{conj:tsang}
 There exists a constant $c > 0$ such that for every Boolean function
 $f$ with Fourier sparsity $s$, there exists an affine subspace of
 codimension $O\left(\log^c s\right)$ on which $f$ is constant.
\end{conjecture} 
Even before, Montanaro and Osborne~\cite{DBLP:journals/corr/abs-0909-3392} had conjectured that the parity decision tree complexity of a function is polylogarithmic in its Fourier sparsity. Tsang \etal\ showed that this seemingly stronger conjecture follows from \expref{Conjecture}{conj:tsang}.

Tsang~\etal~proved the above conjecture for certain classes of
functions, which include functions with constant $\mathbb{F}_2$
degree, and proved \lref[Lemma]{lem:set} for general functions. It follows from %
subsequent %
result by Hatami, Hosseini and Lovett~\cite{DBLP:conf/focs/HatamiHL16}
that \expref{Conjecture}{conj:tsang} is equivalent to the log-rank conjecture for XOR functions (and also to the conjecture by Montanaro and Osborne).


We remark that the bound of \lref[Theorem]{thm:th1} is near-optimal; so any significant improvement to it assuming \expref{Conjecture}{conj:tsang} would be a contradiction, and hence would serve to refute \lref[Conjecture]{conj:tsang}. We note that with our proof technique and analysis, any improvement
to \lref[Lemma]{lem:set} (in particular a positive resolution of
\lref[Conjecture]{conj:tsang}), does not yield a better than
logarithmic improvement to \lref[Theorem]{thm:th1}. Our result, thus, does not seem to throw any light on the truth of \lref[Conjecture]{conj:tsang}. For further
discussion on this topic, the reader is referred to \lref[Section]{themain}.

\section{Preliminaries}\label{sec:prelim}

Let $f:\F_2^n \rightarrow \{1,-1\}$ be a Boolean function. For $\gamma=(\gamma_1, \ldots, \gamma_n) , x=(x_1, \ldots, x_n) \in \F_2^n$, let $\chi_\gamma(x):=(-1)^{\sum_{i=1}^n \gamma_ix_i}$ be the \emph{character} associated with $\gamma$. It is well known that
every Boolean function $f(x)$
can be uniquely written as
\begin{align}f(x)=\sum_{\gamma \in \mathbb{F}_2^n}\widehat{f}(\gamma) \chi_\gamma(x).\label{FT}\end{align}
The right-hand side of \expeqref{Equation}{FT} is called the \emph{Fourier expansion} of $f$, and the real coefficients $\widehat{f}(\gamma)$ are called the \emph{Fourier coefficients}.

We review some standard definitions and facts about the Fourier
coefficients.
\begin{definition}
Let $f(x)=\sum_{\gamma \in \mathbb{F}_2^n}\widehat{f}(\gamma) \chi_\gamma(x)$ be a Boolean function, and $p \geq 1$. The $p$-th spectral norm $\widehat{\|}f\widehat{\|}_p$ of $f$ is defined as %
\[\widehat{\|}f\widehat{\|}_p=\left(\sum_{\gamma \in \mathbb{F}_2^n} \left|\widehat{f}(\gamma)\right|^p\right)^{1/p}.\]

\end{definition}

\begin{lemma}[A special case of Parseval's identity]
\label{parse}
 For a Boolean function $f$, $\widehat{\|}f\widehat{\|}_2 = 1$.
\end{lemma}
The 1st spectral norm of a Boolean function can be bounded in terms of sparsity as follows.
\begin{clm}
\label{claim:lonespar}
For a Boolean function $f$ with Fourier sparsity $s$, $\widehat{\|}f\widehat{\|}_1 \leq \sqrt{s}$.
\end{clm}
\begin{proof}
\begin{align*}
\widehat{\|}f\widehat{\|}_1  \leq \widehat{\|}f\widehat{\|}_2 \cdot \sqrt{s}   = \sqrt{s}.
\end{align*}
The first inequality is the Cauchy-Schwarz inequality while the
second equality follows from \lref[Lemma]{parse}.
\end{proof}

In order to prove our results, we shall use the following version of the Uncertainty Principle. See, e.g., \cite{DBLP:journals/cjtcs/GurT13} for a proof.
\begin{theorem}[Uncertainty Principle]
\label{thm:up}
Let $p: \mathbb{R}^n \rightarrow \mathbb{R}$ be a real
%
multilinear   %
non-zero
 $n$-variate polynomial with sparsity $s$ (\ie, it has $s$ monomials with non-zero coefficients). Let $U_n$ denote the 
uniform distribution on $\{1,-1\}^n$. Then
\[\Pr_{x \sim U_n}[p(x) \neq 0] \geq \frac{1}{s}.\]
\end{theorem}
As
%
mentioned  %
in the
%
Introduction,    %
we need the following
%
%
%
%
result by Tsang~\etal~\cite[Lemma 28]{tsang}.
\begin{theorem}[Tsang~\etal]
\label{sphilka}
Let $f:\mathbb{F}_2^n \rightarrow \{1,-1\}$ be such that
$\widehat{\|}f\widehat{\|}_1=A$. Then there is an affine subspace
$V$ of $\mathbb{F}_2^n$ of codimension $O(A)$ such
that $f$ is constant on $V$.
\end{theorem}

\lref[Lemma]{lem:set} is a simple corollary of this theorem via \lref[Claim]{claim:lonespar}.

In our proof we crucially use the observation made by Gopalan~\etal~\cite{DBLP:journals/siamcomp/GopalanOSSW11} about the equivalence of the non-adaptive parity decision tree complexity of a function (defined below) and its Fourier dimension. We state it in \expref{Proposition}{prop:basic}.
\begin{definition}[Non-adaptive parity decision tree complexity]
Let $f$ be a Boolean function. The \emph{non-adaptive parity decision
  tree complexity} of $f$ (denoted by $\text{NADT}_\oplus(f)$) is defined as the minimum integer $t$ such that there
exist $\gamma_1, \ldots, \gamma_t \in \mathbb{F}_2^n$ such that $f$ is determined by the evaluation of the characters $\chi_{\gamma_1}, \ldots, \chi_{\gamma_t}$.
\end{definition}
\begin{proposition}[Gopalan \etal~\cite{DBLP:journals/siamcomp/GopalanOSSW11}]
\label{prop:basic}
 For every Boolean function $f$, $\mbox{NADT}_\oplus(f) = \dimn(\widehat{f})$.
\end{proposition}
\begin{comment}
\begin{proof}
If the outputs of a basis of span of $\Supp(\widehat{f})$ are specified, then that clearly specifies the outputs of all characters in $\Supp(\widehat{f})$, and hence it specifies the output of the function. Thus $\mbox{NADT}_\oplus(f) \leq \dimn(\widehat{f})$. \\ \\
Now, Let $\mbox{NADT}_\oplus(f)=t$. Let the outputs of $\chi_{\gamma_1}, \ldots, \chi_{\gamma_t}$ specify the output of $f$. $\gamma_i$'s are linearly independent as vectors in $\mathbb{F}_2^n$ (else a smaller number of characters would decide the output of
$f$).
Arbitrarily extend $\gamma_1, \ldots, \gamma_t$ to a basis $\gamma_1, \ldots, \gamma_n$ of $\mathbb{F}_2^n$.
For $x \in \mathbb{F}_2^n$, let $L(x)=(\gamma_1(x), \ldots, \gamma_n(x))$. $L$ is easily seen to be an invertible linear transformation from $\mathbb{F}_2^n$
onto itself. Now, $\forall x \in \mathbb{F}_2^n, \forall i \in [n], \gamma_i(x)=(L(x))_i$. Replacing $x$ by $L^{-1}(x)$ we have $\gamma_i(L^{-1}(x))=x_i$. Now consider the Boolean function
$g(x)=f(L^{-1}(x))=\sum_{\gamma \in \widehat{\mathbb{F}_2^n}}\widehat{f}(\gamma) (-1)^{\gamma(L^{-1}(x))}$. Clearly $\dimn(\widehat{g})=\dimn(\widehat{f})$ as $L$ is a full-rank linear transformation. Also, $g$ is completely specified by the outputs
of $\gamma_i(L^{-1}(x))$'s for $i \in [t]$. Since $\gamma_i(L^{-1}(x))=x_i$, we have that $g$ is determined by $x_1, \ldots, x_t$. Thus all the monomials in $\Supp(\widehat{g})$ contain only the variables $x_1, \ldots, x_t$. Thus $\dimn(\widehat{f}) = \dimn(\widehat{g})
\leq t = \mbox{NADT}_\oplus(f)$. \\ \\
The proposition follows by combining the two inequalities.
\end{proof}
\end{comment}
%
%

\section{Upper bound on parity decision tree complexity} %
\label{themain}
In this section, we %
prove an upper bound on the non-adaptive parity decision tree
complexity of a Boolean function $f$ in terms of its Fourier sparsity $s$. For $B \subseteq \F_2^n$, let $f \mid_B$ denote the restriction of the function $f$ to the set $B$. We will often identify vectors $\gamma \in \F_2^n$ with the linear function that maps a vector $x \in \F_2^n$ to $\gamma(x):=\sum_{i=1}^n \gamma_ix_i \bmod 2 \in \F_2$. We first formally present a procedure {\sc naDT} that constructs a non-adaptive
parity decision tree of $f$. We then provide a description of the procedure in words (in particular, the roles of the various variables used) and a formal analysis of the same, leading up to the proof of \lref[Theorem]{thm:th1}.

\newcommand{\nadtproc}{\text{\sc naDT}}
\begin{itemize}
\item[] $\nadtproc(f)$
\item[] Input: Boolean function $f: \F_2^n \to \{1,-1\}$;
\item[] Output: A set $\Gamma$ of parities, whose evaluations determine the value of $f$;
\begin{enumerate}
\item Set $\Gamma \gets \varnothing$, $\mathcal{S} \gets\Supp(\widehat{f})$
  and 
  $\mathcal{F} \gets \{f\}$.
\item\label{step:while} While $\mathcal{S} \neq \{\emptyset\}$, do
\begin{enumerate}
\item\label{step:chooseg} Let $g$ be a function in $\mathcal{F}$ with the largest Fourier
  sparsity. Let $A$ be a largest
  affine subspace on which $g$ is constant (breaking ties arbitrarily), with codimension $n_g$.  Let $\gamma_1,\ldots,\gamma_{n_g} \in \F_2^n$ and
  $b_1, \ldots, b_{n_g} \in \mathbb{F}_2$ be such that
  \[
    A=\{x \in \F_2^n :
    \gamma_1(x)=b_1, \ldots, \gamma_{n_g}(x) = b_{n_g}\}.
  \]
\item Set $\Gamma \gets \Gamma \cup \{\gamma_1,\ldots,
  \gamma_{n_g}\}$.
\item For each $b=(b_\gamma)_{\gamma \in \Gamma} \in
  \mathbb{F}_2^{|\Gamma|}$, let $V_b$ be the affine subspace $\{x \in \F_2^n:
  \forall \gamma \in \Gamma, \gamma(x) = b_\gamma\}$.  Set
  \[
    \mathcal{F} \gets \bigcup_{b \in \mathbb{F}_2^{|\Gamma|}} \{f |_{V_b}\}.
  \]
\item $\mathcal{S} \gets \bigcup_{h \in \mathcal{F}} \Supp(\widehat{h})$.
\end{enumerate}
\item Return $\Gamma$.
\end{enumerate}
\end{itemize}
\noindent   %
Notation.  %
After each iteration of the while loop in the procedure,
$\Gamma$ is the set of parities that have been queried so far,
$\mathcal{F}$ is the set of all restrictions of $f$ to the affine
subspaces obtained by different assignments to parities in $\Gamma$,
and $\mathcal{S}$ is the union of the Fourier supports of functions in
$\mathcal{F}$. Throughout this section, $i$ will stand for the index of an arbitrary iteration
of {\sc naDT}. Let  $\Gamma^{(i)}, \mathcal{F}^{(i)}$ and
$\mathcal{S}^{(i)}$ denote $\Gamma, \mathcal{F}$ and $\mathcal{S}$
respectively at the end of the $i$-th iteration of the while loop. Let  $\Gamma^{(0)}=\varnothing$,  
$\mathcal{F}^{(0)}=\{f\}$ and $\mathcal{S}^{(0)}=\Supp(\widehat{f})$.


For each $i$, let $b=(b_\gamma)_{\gamma \in \Gamma^{(i)}} \in
\mathbb{F}_2^{|\Gamma^{(i)}|}$ and let $V_{b}$ be the affine subspace
defined by the linear constraints $\{\gamma(x)=b_\gamma: \gamma \in
\Gamma^{(i)}\}$.  In $V_b$, more than one linear function of the
original space may get identified,
%
namely,  %
the restrictions of these linear functions to the
sub-domain  %
subdomain   %
$V_b$ are either the same function or negations of each
  other. More specifically, $\delta_1$ and $\delta_2$ get
identified in $V_b$ if and only if $\delta_1 + \delta_2 \in
\vspan\Gamma^{(i)}$ (\ie, %
they belong to the same coset of the subspace $\vspan\Gamma^{(i)}$). This equivalence relation
induces a partition on $\Supp(\widehat{f})$
into equivalence classes. Note that this partition is determined by $\Gamma^{(i)}$ alone.
Further, for each equivalence class, for every $b \in
\mathbb{F}_2^{|\Gamma^{(i)}|}$, the linear functions belonging to
that class get identified with one another in $V_b$. 

Let $\ell^{(i)}$ denote the number of such equivalence classes.\footnote{$\ell^{(i)}$ is in fact the number of cosets of the subspace
$\vspan\Gamma^{(i)}$ with which $\Supp(\widehat{f})$ has non-empty
intersection.} For all $j \in [\ell^{(i)}]$, let $\beta_j^{(i)}$ be an arbitrarily picked
representative element in the $j$-th equivalence class. Let
\[
  \beta_j^{(i)} + \alpha_{j,1}^{(i)},
  \ldots, \beta_j^{(i)} + \alpha_{j,{k_j^{(i)}}}^{(i)}
\]
be the $k_j^{(i)}$ elements in the $j$-th equivalence class, where
$\alpha_{j,1}^{(i)}, \ldots, \alpha_{j,{k_j^{(i)}}}^{(i)}$ are in
$\vspan\Gamma^{(i)}$. Now define functions
\[
  P_j^{(i)}(x) := \displaystyle\sum_{l=1}^{k_j^{(i)}}
  \widehat{f}\left(\beta_j^{(i)}+\alpha_{j,l}^{(i)}\right)\chi_{\alpha_{j,l}^{(i)}}(x).
\]
Note that the functions $P_j^{(i)}$ are non-zero.  As we will soon
see, it will be helpful to think of the $P_j^{(i)}$'s as
%
multilinear   %
polynomials in variables $y_i=(-1)^{x_i}$.

Given this notation, we can then write the Fourier expansion of $f$ in
the following form:
\[f(x)=\displaystyle\sum_{j=1}^{\ell^{(i)}}P_j^{(i)}(x)\chi_{\beta_j^{(i)}}(x).\]

Note that the value of each $P_j^{(i)}$ is fixed once the evaluation of each linear form $\gamma \in \Gamma^{(i)}$ is specified. In other words, each $P_j^{(i)}$ is a
constant function on each $V_b$.

\begin{observation}
\label{obs:sum}
$\displaystyle\sum_{j=1}^{\ell^{(i)}}k_j^{(i)}=s$.
\end{observation}
\begin{proof}
The size of the Fourier support of $f$ is the sum of the sizes of its equivalence classes defined above.
\end{proof}
\begin{proposition}
\label{prop:lbound}
$|\mathcal{S}^{(i)}| = \ell^{(i)}$.
\end{proposition}
\begin{proof}
 Clearly $|\mathcal{S}^{(i)}| \leq \ell^{(i)}$, as each term in the Fourier expansion of $f \mid_{V_b}$ corresponds to a distinct equivalence class (and this correspondence is independent of $b$). Now, for $j\in[\ell^{(i)}]$, since $P_j^{(i)}$ is a non-zero function, there exists an assignment $b$ to the parities in $\Gamma^{(i)}$ on which $P_j^{(i)}$ evaluates to
a non-zero value. Thus the coefficient of $\beta_j^{(i)}$ is non-zero in the restriction of $f$ to the affine subspace obtained by assigning $b$ to the parities in $\Gamma^{(i)}$.
Thus for all $j\in [\ell^{(i)}]$, $\beta_j^{(i)} \in \mathcal{S}^{(i)}$ which, together with $|\mathcal{S}^{(i)}| \leq \ell^{(i)}$, implies $|\mathcal{S}^{(i)}| = \ell^{(i)}$.
\end{proof}
We now argue that after every iteration of the while loop, there
exists a function $h \in \mathcal{F}^{(i)}$ which has large Fourier support.
\begin{lemma}
\label{lem:support}
After the $i$-th iteration, there exists an $h \in \mathcal{F}^{(i)}$ such
that $|\Supp(\widehat{h})|$ is at least $\left({\ell^{(i)}}\right)^2/{s}$.
\end{lemma}
\begin{proof}
Consider any function $f|_{V_b} \in \mathcal{F}^{(i)}$. The Fourier
decomposition of $f|_{V_b}$ is given by
\[
  f|_{V_b} =
  \sum_{j=1}^{\ell^{(i)}}P_j^{(i)}(b)\cdot \chi_{\beta_j^{(i)}}(x).
\]
Thus, $|\Supp(\widehat{f|_{V_b}})|$ is exactly the number of functions
$P_j^{(i)}, j\in [\ell^{(i)}]$ such that $P_j^{(i)}(b)$ is
non-zero. We analyze this quantity as follows. Pick a
$b \in \mathbb{F}_2^{|\Gamma^{(i)}|}$ uniformly at random. Now, note
that for each $j \in [\ell^{(i)}]$, $P_j^{(i)}$ is a non-zero
%
multilinear   %
polynomial
in variables $\{(-1)^{b_\gamma}: \gamma \in \Gamma^{(i)}\}$. Thus by \lref[Theorem]{thm:up},
$$\Pr_b[P_j^{(i)}(b) \neq 0] \geq \frac{1}{k_j^{(i)}}.$$ Thus,
\begin{align*}
\mathbb{E}_b\left[|\Supp(\widehat{f|_{V_b}})|\right] &\geq \sum_{j=1}^{\ell^{(i)}}
\frac{1}{k_j^{(i)}} \\
&\geq \ell^{(i)} \cdot \frac{1}{\left(\sum_{j=1}^{\ell^{(i)}}
    k_j^{(i)}\right)/\ell^{(i)}}  && \text{[By Jensen's inequality applied to $1/x$]}\\
&= \frac{\left(\ell^{(i)}\right)^2}{s} && \text{[By
  \lref[Observation]{obs:sum}]}.
\end{align*}
Hence, there exists an $h \in \mathcal{F}^{(i)}$ such
that $|\Supp(\widehat{h})|$ is at least $\left({\ell^{(i)}}\right)^2/{s}$.
\end{proof}
Let $g^{(i)}$ be the function chosen in \lref[step]{step:chooseg} of the
$i$-th iteration of $\nadtproc$. Let $\spa(\widehat{g^{(i)}}) = s^{(i)}$, and $\Delta \ell^{(i)}= \ell^{(i-1)} - \ell^{(i)}$ for $i \geq 1$. The next lemma proves that, if a function with large Fourier support is picked in \lref[step]{step:chooseg}, then that leads to a large reduction in
the size of $\mathcal{S}$.
\begin{lemma}
\label{lem:reduction}
Assume that neither $f$ nor $1-f$ is a character, and $\nadtproc(f)$ runs for $t$ iterations. Then for all $i\in [t]$, $\Delta \ell^{(i)} \geq {s^{(i)}}/{2}$.
\end{lemma}
\begin{proof}
  Let $\gamma_1,\ldots,\gamma_{n_{g^{(i)}}}$ be the parities queried
  in iteration $i$.  Hence there is a
  $b=(b_1,\ldots,b_{n_{g^{(i)}}}) \in
  \left(\mathbb{F}_2\right)^{n_{g^{(i)}}}$ such that $g^{(i)}$ is
  constant on the affine subspace $V_b$ obtained by setting $\gamma_j$
  to $b_j$ for all $j \in [n_{g^{(i)}}]$. Now, consider the following
  cases.
  \begin{description}
 %
   \item [Case 1. $s^{(i)} \geq 2$.]  %
  Since $g^{(i)}$ is
  constant on $V_b$, each Fourier coefficient of $g^{(i)}\mid_{V_b}$ must be either $0$ or $\pm 1$. This is possible only if for every $\alpha_1 \in \Supp(\widehat{g^{(i)}})$ there is another $\alpha_2 \in \Supp(\widehat{g^{(i)}})$ such that $\alpha_1+\alpha_2 \in \vspan \{\gamma_1, \ldots, \gamma_{n_{g^{(i)}}}\}$, \ie, they get identified. This implies that for every
  \[
    b' =(b')_j \in
    \left(\mathbb{F}_2\right)^{n_{g^{(i)}}},
  \]
  in the affine space $V_{b'}$ obtained by restricting each $\gamma_j$
  to $b'_j$, every parity in $\Supp(\widehat{g^{(i)}})$ is identified
  with some other parity in $\Supp(\widehat{g^{(i)}})$. Since
  $\Supp(\widehat{g^{(i)}}) \subseteq \mathcal{S}^{(i-1)}$, it follows
  that $|\mathcal{S}^{(i)}|$ is at least
  $|\Supp(\widehat{g^{(i)}})|/{2}$ less than
  $|\mathcal{S}^{(i-1)}|$. By \lref[Proposition]{prop:lbound} this
  implies $\Delta \ell^{(i)} \geq {s^{(i)}}/{2}$.
%
 \item [Case 2. $s^{(i)}=1$]   %
(\ie, %
$g^{(i)}$ is either a parity function or the complement of a parity function). Since neither $f$ nor $1-f$ is a character, we have that $s^{(1)}=\spa(\widehat{f}) \geq 2$. Hence we conclude that $i \geq 2$. Note that for $i \geq 1$, $\emptyset \in \mathcal{S}^{(i)}$. Hence for $i \geq 2$, $\Delta \ell^{(i)}=1 \geq {s^{(i)}}/{2}$.
  \end{description}
The lemma follows from the two cases.
\end{proof}
Now we are ready to prove \lref[Theorem]{thm:th1}.
\begin{proof}[Proof of Theorem \lref{thm:th1}]
If either $f$ or $1-f$ is a character, the theorem follows immediately. Hence assume that neither $f$ nor $1-f$ is a character.

We obtain a non-adaptive parity decision tree for $f$ by running $\nadtproc$. Assume that the while loop runs for $t$ iterations.
Let the number of queries made in \lref[step]{step:chooseg} in the $i$-th iteration
of the procedure be $q^{(i)}$. By \lref[Lemma]{lem:set}, $q^{(i)} = O(\sqrt {s^{(i)}})$. By \lref[Lemma]{lem:reduction}, $\Delta \ell^{(i)} \geq {s^{(i)}}/{2}$.
Hence,
\[
  \frac {q^{(i)}}{\Delta \ell^{(i)}}=\frac{1}{\Omega(\sqrt{s^{(i)}})}.
\]
{}From \lref[Lemma]{lem:support} we have
$s^{(i)} \geq \left(\ell^{(i-1)}\right)^2/s$. Hence
$q^{(i)}=\sqrt{s} \cdot O\left(\Delta \ell^{(i)}/\ell^{(i-1)}
\right)$. Thus the total number of queries made within the while loop
of the procedure is
\begin{align}
\label{count}
 \displaystyle\sum_{i=1}^t q^{(i)} &= \sqrt{s} \cdot \displaystyle\sum_{i=1}^t O\left(\frac{\Delta \ell^{(i)}}{\ell^{(i-1)}}\right) \nonumber 
                                           = \sqrt{s} \cdot \displaystyle\sum_{i=1}^t O\left(\frac{\Delta \ell^{(i)}}{s - \displaystyle\sum_{j=1}^{i-1} \Delta \ell^{(j)}}\right) \nonumber \\
					  %
					  & \leq \sqrt{s} \cdot \displaystyle\sum_{i=1}^t O \left(  \frac{1}{s - \displaystyle\sum_{j=1}^{i-1} \Delta \ell^{(j)}} + \frac{1}{s - \displaystyle\sum_{j=1}^{i-1} \Delta \ell^{(j)}-1} 
+\cdots+  \frac{1}{s - \displaystyle\sum_{j=1}^{i-1} \Delta \ell^{(j)}-(\Delta \ell^{(i)}-1)} \right) \nonumber \\
					  &\leq \sqrt{s} \cdot \displaystyle O\left(\sum_{k=1}^s \frac{1}{k}\right) = O(\sqrt{s} \log s). \nonumber
\end{align}
{}From \lref[Proposition]{prop:basic} it follows that
$\dimn(\widehat{f}) = O(\sqrt{s} \log s)$.
\end{proof}

\paragraph{Discussion.}  %
As mentioned in the introduction, a potential approach towards disproving \lref[Conjecture]{conj:tsang} is to assume it to be true, and prove that it implies a $o(\sqrt{s})$ upper bound on Fourier dimension.
This will refute the conjecture, since, for the address function (see \lref[Section]{intro}),
\[
  \dimn(\widehat{\Add_s}) = \Theta\left(\sqrt{\spa(\widehat{\Add_s)}}\right).
\]
However, we cannot disprove the conjecture by an analysis of the
$\nadtproc$, assuming the conjecture. To see this let us consider the
execution of {\sc naDT} on the address function. Recall that
$\Add_s(x,y_1,y_2,\dots,y_{\sqrt{s}}) = y_x, x \in\{0,1\}^{(1/2)\log
  s}, y_i \in \{0,1\}$. One easily sees that a largest affine subspace
$V$ on which the function is constant is the one defined by the
constraints $x=x'$, $y_{x'}=b$ where
$x' \in \{0,1\}^{({1}/{2}) \log s}$ and $b \in \{0,1\}$. The
function takes the value $b$ everywhere in $V$. Also, if the
%
address     %
bits $x$ and the bit $y_{x'}$ are set to other values than $x'$ and $b$,
then the restricted functions in the respective affine subspaces are all
constants (if $x$ is set to $x'$) or dictators (on
%
the addressed  %
bits). The subsequent steps
%
query  %
different dictators on the
%
addressed   %
bits.

The address function clearly satisfies \lref[Conjecture]{conj:tsang}, and all the intermediate functions that are given rise to by $\nadtproc$ are dictators, which also trivially satisfy the conjecture.
Thus this rules out the possibility of refuting \lref[Conjecture]{conj:tsang} by analyzing $\nadtproc$ assuming the conjecture. We note, however, that if we assume the conjecture, then we
can improve the upper bound by a factor of $O(\log s)$, to the optimal $O(\sqrt{s})$.
\section*{Acknowledgements}
I thank Avishay Tal for observing that an earlier analysis of $\nadtproc$ which proved a $O(s^{2/3})$ upper bound can be tightened to obtain the $O(\sqrt{s} \log s)$ upper bound, and bringing it to my notice. I thank the anonymous referees for pointing out that \lref[Theorem]{thm:th1} implies
the  %
existence of simultaneous communication protocols for XOR functions, for suggesting simplifications of some proofs, and for suggestions that helped to %
improve the presentation. I thank Arkadev Chattopadhyay and Prahladh Harsha for many helpful discussions. I thank Prahladh Harsha for his help in
%
significantly improving the presentation of this article.  %

%% Theory of Computing
%% File: v015a011.bbl
%% Publication date: October 27, 2019
%% Authors: Swagato Sanyal
% Define empty bibhead if not already defined
\providecommand{\bibhead}[1]{}
% Define tocrefpdfbookmark if not already defined
\expandafter\ifx\csname pdfbookmark\endcsname\relax%
  \providecommand{\tocrefpdfbookmark}{}
\else\providecommand{\tocrefpdfbookmark}{%
   \hypertarget{tocreferences}{}%
   \pdfbookmark[1]{References}{tocreferences}}%
\fi

\tocrefpdfbookmark
\begin{thebibliography}{10}

\bibitem{DBLP:journals/corr/abs-1810-00481}\bibhead{DBLP:journals/corr/abs-1810-00481}
{\sc Srinivasan Arunachalam, Sourav Chakraborty, Troy Lee, Manaswi Paraashar,
  and Ronald de~Wolf}: Two new results about quantum exact learning.
\newblock In {\em Proc. 46th Internat. Colloq. on Automata, Languages and
  Programming (ICALP'19)}, pp. 16:1--16:15. Schloss Dagstuhl--Leibniz-Zentrum
  fuer Informatik, 2019.
\newblock [\epfmtdoi{10.4230/LIPIcs.ICALP.2019.16}, \epfmt{arxiv}{1810.00481}]

\bibitem{DBLP:journals/eccc/ChattopadhyayMS18}\bibhead{DBLP:journals/eccc/ChattopadhyayMS18}
{\sc Arkadev Chattopadhyay, Nikhil~S. Mande, and Suhail Sherif}: The
  log-approximate-rank conjecture is false.
\newblock In {\em Proc. 50th STOC}, pp. 42--53. ACM Press, 2019.
\newblock [\epfmtdoi{10.1145/3313276.3316353}]

\bibitem{alex}\bibhead{alex}
{\sc Dmitry Gavinsky, Naomi Kirshner, Alex Samorodnitsky, and Ronald de~Wolf}:
  Private communication, 2014.

\bibitem{DBLP:journals/siamcomp/GopalanOSSW11}\bibhead{DBLP:journals/siamcomp/GopalanOSSW11}
{\sc Parikshit Gopalan, Ryan O'Donnell, Rocco~A. Servedio, Amir Shpilka, and
  Karl Wimmer}: Testing {F}ourier dimensionality and sparsity.
\newblock {\em SIAM J. Comput.}, 40(4):1075--1100, 2011.
\newblock Preliminary version in
  \href{https://link.springer.com/chapter/10.1007/978-3-642-02927-1_42}{ICALP'09}.
\newblock [\epfmtdoi{10.1137/100785429}]

\bibitem{DBLP:journals/cjtcs/GurT13}\bibhead{DBLP:journals/cjtcs/GurT13}
{\sc Tom Gur and Omer Tamuz}: Testing {B}ooleanity and the uncertainty
  principle.
\newblock {\em Chicago J. Theoret. Computer Sci.}, 2013(14), 2013.
\newblock [\epfmtdoi{10.4086/cjtcs.2013.014}, \epfmt{arxiv}{1204.0944}]

\bibitem{DBLP:conf/focs/HatamiHL16}\bibhead{DBLP:conf/focs/HatamiHL16}
{\sc Hamed Hatami, Kaave Hosseini, and Shachar Lovett}: Structure of protocols
  for {XOR} functions.
\newblock {\em SIAM J. Comput.}, 47(1):208--217, 2018.
\newblock Preliminary version in
  \href{https://ieeexplore.ieee.org/document/7782941}{FOCS'16}.
\newblock [\epfmtdoi{10.1137/17M1136869}]

\bibitem{DBLP:journals/toct/HavivR16}\bibhead{DBLP:journals/toct/HavivR16}
{\sc Ishay Haviv and Oded Regev}: The list-decoding size of {F}ourier-sparse
  {B}oolean functions.
\newblock {\em ACM Trans. Comput. Theory}, 8(3):10:1--10:14, 2016.
\newblock Preliminary version in
  \href{http://drops.dagstuhl.de/opus/volltexte/2015/5060/}{CCC'15}.
\newblock [\epfmtdoi{10.1145/2898439}, \epfmt{arxiv}{1504.01649}]

\bibitem{DBLP:conf/coco/KannanMSY18}\bibhead{DBLP:conf/coco/KannanMSY18}
{\sc Sampath Kannan, Elchanan Mossel, Swagato Sanyal, and Grigory
  Yaroslavtsev}: Linear sketching over $\mathbb{F}_2$.
\newblock In {\em Proc. 33rd Computational Complexity Conf. (CCC'18)}, pp.
  8:1--8:37. Schloss Dagstuhl--Leibniz-Zentrum fuer Informatik, 2018.
\newblock [\epfmtdoi{10.4230/LIPIcs.CCC.2018.8}, \epfmt{arxiv}{1611.01879}]

\bibitem{khalyavin}\bibhead{khalyavin}
{\sc Andrei~Vyacheslavovich Khalyavin, Mikhail~Sergeevich Lobanov, and
  Yuriy~Valerievich Tarannikov}: On plateaued {B}oolean functions with the same
  spectrum support.
\newblock {\em Sib. \`Elektron. Mat. Izv.}, 13:1346--1368, 2016.
\newblock [\epfmtdoi{10.17377/semi.2016.13.105}]

\bibitem{lovett}\bibhead{lovett}
{\sc Shachar Lovett}: Communication is bounded by root of rank.
\newblock {\em J. {ACM}}, 63(1):1:1--1:9, 2016.
\newblock Preliminary version in
  \href{http://doi.org/10.1145/2591796.2591799}{STOC'14}.
\newblock [\epfmtdoi{10.1145/2724704}, \epfmt{arxiv}{1306.1877}]

\bibitem{DBLP:journals/corr/abs-0909-3392}\bibhead{DBLP:journals/corr/abs-0909-3392}
{\sc Ashley Montanaro and Tobias Osborne}: On the communication complexity of
  {XOR} functions, 2009.
\newblock [\epfmt{arxiv}{0909.3392}]

\bibitem{conf-version}\bibhead{conf-version}
{\sc Swagato Sanyal}: Near-optimal upper bound on {F}ourier dimension of
  {B}oolean functions in terms of {F}ourier sparsity.
\newblock In {\em Proc. 42nd Internat. Colloq. on Automata, Languages and
  Programming (ICALP'15)}, pp. 1035--1045. Springer, 2015.
\newblock [\epfmtdoi{10.1007/978-3-662-47672-7\_84}]

\bibitem{DBLP:conf/innovations/ShpilkaTV14}\bibhead{DBLP:conf/innovations/ShpilkaTV14}
{\sc Amir Shpilka, Avishay Tal, and Ben lee Volk}: On the structure of
  {B}oolean functions with small spectral norm.
\newblock {\em Comput. Complexity}, 26(1):229--273, 2017.
\newblock Preliminary version in
  \href{https://dl.acm.org/citation.cfm?doid=2554797.2554803}{ITCS'14}.
\newblock [\epfmtdoi{10.1145/2591796.2591799}, \epfmt{arxiv}{1304.0371}]

\bibitem{tsang}\bibhead{tsang}
{\sc Hing~Yin Tsang, Chung~Hoi Wong, Ning Xie, and Shengyu Zhang}: Fourier
  sparsity, spectral norm, and the {L}og-rank {C}onjecture.
\newblock In {\em Proc. 54th FOCS}, pp. 658--667. IEEE Comp. Soc. Press, 2013.
\newblock [\epfmtdoi{10.1109/FOCS.2013.76}, \epfmt{arxiv}{1304.1245}]

\end{thebibliography}

\end{document}
